\chapter{Sustancias que no definen el tipo de neurona}
\label{sec:othermediators}

En el capítulo~\ref{sec:neurontypes} clasificamos las neuronas por su neurotransmisor principal y obtuvimos diez tipos (tabla~\ref{tab:types}). Es hora de añadir complicaciones. Una neurona no libera solo su neurotransmisor principal, y en el cerebro, además de los neurotransmisores principales, actúan otras sustancias. Cambian la manera en que la red transmite y almacena las señales.

Además del bus de direcciones y del bus de difusión (sección~\ref{sec:buses}) hay un tercero: el \emph{canal retrógrado}. Varias sustancias no las libera el emisor, sino el \emph{receptor}, y viajan por la hendidura hacia atrás, hasta la salida del emisor, y cambian cuánto neurotransmisor liberará este la próxima vez. En las redes de comunicación se parece a un acuse de recibo que el emisor usa para ajustar su transmisión. Es precisamente este canal el que usan las neuronas cuando cambian los pesos de sus conexiones: por él, la salida del emisor se entera de la actividad del receptor, es decir, del segundo factor de las reglas de aprendizaje del capítulo~\ref{sec:dofa}.

La lista completa de neurotransmisores conocidos es larga: más de cien sustancias, la mayoría cadenas cortas de proteína, los neuropéptidos. Pero todas caben en unas pocas clases, y dentro de cada clase las sustancias se comportan de forma parecida. Las sustancias que nunca son el neurotransmisor principal están reunidas en la tabla~\ref{tab:othermediators}, con las mismas tres preguntas que haría un ingeniero: por qué bus va la señal, si actúa rápido y cuál suele ser su signo. No definen el tipo de neurona: se liberan junto con el neurotransmisor principal, o no las liberan neuronas, o van en sentido contrario.

\begin{table}[t]
\centering
\footnotesize
\setlength{\tabcolsep}{4pt}
\renewcommand{\arraystretch}{1.12}
\begin{tabular}{>{\raggedright}p{2.25cm}>{\raggedright}p{2.8cm}>{\raggedright}p{1.8cm}>{\raggedright}p{1.5cm}>{\centering}p{0.8cm}>{\raggedright\arraybackslash}p{4.3cm}}
\toprule
Clase & Sustancia & Bus & Velocidad & Signo & Papel en el cálculo \\
\midrule
Aminoácidos & D-serina & local & lenta & Y & segunda llave del receptor de glutamato: condición “Y” \\
\midrule
Monoaminas & aminas traza & difusión & lenta & $\pm$ & ajuste fino de los canales monoaminérgicos \\
\midrule
\multirow{2}{2.25cm}{Purinas}
 & ATP & de direcciones & rápida y lenta & $+$ & cotransmisor; señal entre neuronas y glía \\
 & adenosina & de volumen & lenta & $-$ & se acumula con la actividad: contador de carga~\cite{Dunwiddie2001} \\
\midrule
\multirow{8}{2.25cm}{Neuropéptidos (más de cien)}
 & encefalinas, endorfinas, dinorfina & de volumen & lenta & $-$ & supresión de la transmisión \\
 & sustancia P & de volumen & lenta & $+$ & amplificación lenta \\
 & neuropéptido Y & de volumen & lenta & $-$ & inhibición lenta \\
 & somatostatina & de volumen & lenta & $-$ & inhibición lenta, acompaña al GABA \\
 & péptido intestinal vasoactivo (VIP) & de volumen & lenta & $+$ & desinhibición: inhibición de las neuronas inhibidoras \\
 & colecistoquinina & de volumen & lenta & $\pm$ & acompaña al GABA en parte de las neuronas inhibidoras \\
 & orexina & difusión & lenta & $+$ & estabilizador del modo de vigilia \\
 & oxitocina, vasopresina & difusión & lenta & $\pm$ & ajustes globales del comportamiento \\
\midrule
\multirow{2}{2.25cm}{Gases}
 & óxido nítrico NO & retrógrado, de volumen & lenta & $\pm$ & atraviesa las membranas; señal retrógrada al cambiar los pesos~\cite{Garthwaite2008} \\
 & CO, H$_2$S & de volumen & lenta & $\pm$ & papel poco estudiado \\
\midrule
Lípidos & endocannabinoides (anandamida, 2-AG) & retrógrado & lenta & $-$ & el receptor reduce la liberación del emisor: realimentación sobre el peso~\cite{WilsonNicoll2001} \\
\bottomrule
\end{tabular}
\caption{Sustancias que no definen el tipo de neurona. \emph{Bus}, \emph{velocidad} y \emph{signo}, como en la tabla~\ref{tab:types}; además, bus de volumen: el neurotransmisor se difunde por el espacio extracelular alrededor del punto de liberación; retrógrado: va del receptor de vuelta al emisor; local: actúa cerca del punto de liberación. “Y” es una señal habilitadora: por sí sola no produce corriente, pero sin ella no se activa otra entrada, como la segunda entrada de una compuerta lógica “Y”. Los neuropéptidos casi siempre se liberan junto con el neurotransmisor principal, y sobre todo a alta frecuencia de espigas~\cite{vandenPol2012}.}
\label{tab:othermediators}
\end{table}
